% !TEX root = ../nc-quadric-moduli.tex
% 2026-11-21 / v1.35
\section{相対 Picard 作用による cubic Serre 平方根の構成}
\label{sec:relative-picard-roots}

\paragraph{主張の範囲と依存.}
本稿は標数零における\emph{非線型} geometric cubic 五つ組の普遍族上に、
étale 局所に実際の二周期を構成する。
その平方は四次 tensor の巡回から定まる canonical 四周期に
頂点スカラーを除いて一致する。
全ての恒等式を被約な普遍基底上で証明してから引き戻すので、
非被約な Artin 基底も結論に含む。
原典の既知入力は、Artin--Tate--Van den Bergh の
\emph{相対} Picard 作用と、Van den Bergh の geometric 五つ組の構成である。
相対 tails と有限 foundation の比較、および Serre の双加群比較には、
定理~\ref{thm:relative-cubic-tails-tilting} と定理~\ref{thm:bimodule-line-comparison} を用いる。
本稿は任意の二周期が Serre 平方根であるとは主張しない。
相対化の新規性は未確認である。

\subsection{一次資料の正確な内容}

Van den Bergh \cite[Proposition 5.6.1]{VdB2011} は、
代数閉体 $k$、$\operatorname{char}k\ne2$ 上の全ての quadric
$\mathbb Z$ 代数 $D$ に対して $D\simeq D(2)$ を示す。
全 $\mathbb Z$ 代数の二周期性が主張されている。
その後の Corollary 5.6.2 が偶数部分の通常の次数付き環による表示を与える。

非線型の場合、楕円 helix $(C,(\mathcal L_i)_i)$ に対し
\[
 N=\mathcal L_2\mathcal L_0^{-1}\in\Pic^0(C),\qquad
 \mathcal L_{i+2}\simeq\mathcal L_iN
\]
である。$a\in\Pic^0(C)$、$a^{-2}=N$ を選び、
\cite[Theorem 4.2.3]{VdB2011} の作用
\[
 \eta:\Pic^0(C)\longrightarrow\Aut(C),\qquad
 \eta(a)^*L=L\otimes a^{-\deg L}
 \tag{PR-1}\label{eq:field-eta}
\]
を用いる。$\sigma=\eta(a)$ は各次数二の束 $\mathcal L_i$ を $\mathcal L_{i+2}$ へ移す。
原典の Proposition 6.6 は二周期を標準 hull へ持ち上げるが、
任意の二周期 $\beta$ に対する $\beta^2=\theta$ を主張してはいない。

\begin{proposition}[ATV の相対 Picard 作用]\label{prop:ATV-relative-action}
$C\to S$ を $\PP^2_S$ の三次因子、または
$\PP^1_S\times_S\PP^1_S$ の bidegree $(2,2)$ の因子の平坦族とする。
$P=\Pic^0_{C/S}$ を、各幾何ファイバーの各被約既約成分で次数零の
直線束の類を表す相対 Picard 部分群とする。
$P$ の $C$ への作用があり、その構成は任意の底変更および族の同型に自然である。
作用の符号を \eqref{eq:field-eta} と合わせて $\eta$ と書く。
\end{proposition}

\begin{proof}[相対作用の構成]
\cite[p.~355, Proposition 5.5 and Corollary 5.7]{ArtinTateVanDenBergh1991} は初めから
任意の基底 $S$ 上の因子の族 $C/S$ を扱う。
相対多重次数零の直線束 $a$ と section $p:S\to C$ に対し、
$a\otimes I_p$ が局所的にある一意な section $p'$ の ideal と同型になることを示す。
具体的には
\[
 \pi_*\mathcal Hom(a\otimes I_p,\OO_C)
\]
が階数一で、その局所生成射が単射になり、余核が section の構造層となる。
一意性によりこの操作が貼り合わさり、任意の底変更と両立する群作用となる。
原典の (5.8) の $\eta_a$ はこの作用に $a^{-1}$ を入れた符号である。
したがって smooth elliptic fiber 上では usual translation となり、
Lemma 5.10(iii) が \eqref{eq:field-eta} を与える。
ideal と tensor による一意な構成なので、曲線族の同型への自然性も成り立つ。

Corollary 5.7 は相対的である一方、Lemma 5.10(iii) の記述は体上である。
以下では、この違いを埋める普遍基底上の議論を明記する。
\end{proof}

\subsection{被約な普遍基底で証明してから底変更する}

以下、$k$ は代数閉体、$\operatorname{char}k=0$ とする。
四つの二次元空間 $V_0,\ldots,V_3$ を固定し、
\[
 \mathscr G\subset\PP(V_0\otimes V_1\otimes V_2\otimes V_3)
\]
を geometric 五つ組の開部分とする。
$02\mid13$ flattening の階数一の閉部分を除き、
\[
 S=\mathscr G\setminus\mathscr G_{\mathrm{lin}}
\]
と置く。$S$ は滑らかで被約、既約な有限型 $k$ スキームである。
以下の構成は $S$ の affine open 上で行って貼り合わせてもよい。

普遍直線 $W\subset(V_0\otimes\cdots\otimes V_3)\otimes\OO_S$ を巡回し、
$V_{i+4}=V_i$ として $W_i$ と三次関係 $R_i$ を定める。
これにより普遍 cubic $\mathbb Z$ 代数 $D$ を得る。
$W_i$ は全て同じ基底線 $W$ の巡回像であり、この共通の同定の下で
canonical 四周期
\[
 \theta:D\longrightarrow D(4)
 \tag{PR-2}\label{eq:theta}
\]
は各 $V_i=V_{i+4}$ 上の恒等である。
この記述に $W$ の大域的な生成元は必要ない。
局所生成元 $w_i$ を使えば
\[
 w_i=\sum_a v_a\otimes r_a,
 \qquad w_{i+1}=\sum_a r_a\otimes\theta_i(v_a)
\]
という \cite[\S5.2]{VdB2011} の定義そのものである。

\begin{lemma}[普遍族の次数別平坦性]\label{lem:universal-flatness}
各 $D_{ij}$ は有限局所自由で、標準 cubic 解消は各内部次数で分裂完全である。
その双対複体も AS 双対の所定の次数以外で完全で、残る cohomology は可逆層である。
従ってこれらの提示・解消・双対・関係の交わりは任意の底変更と両立する。
\end{lemma}

\begin{proof}
固定した $i,j$ に対し、$D_{ij}$ は有限個の tensor 積と三次関係からなる
有限自由加群間の行列の余核として提示される。
全ての幾何ファイバーは cubic AS 正則で、$\dim D_{ij}$ は一定である。
従って提示行列のファイバー階数は一定である。
適切な階数の小行列式が単元となる近傍で行列を掃き出すと、
余分な block の全成分は全ての点で零となる。
$S$ は被約なので、その block 自体が零である。
従って余核は局所自由で、提示は局所的に標準形へ分裂する。

標準 cubic 複体の固定内部次数は有限局所自由複体となり、
全幾何ファイバーで完全である。
そのような有限複体は局所的に分裂完全であるから、任意底変更後も完全である。
この解消を使って定義される kernels、images および四次の overlap も
直和因子となり、任意底変更と両立する。
双対複体の固定内部次数も有限局所自由複体である。
各幾何ファイバーの AS 双対により、その各 cohomology の次元は一定であり、
これと各項の階数から各微分の階数も一定となる。
被約性を用いた同じ行列の標準形により、像・核・cohomology は局所直和因子となり、
その形成は任意底変更と両立する。
特に右単純対象 $S_i$ と $P_j=e_jD$ について
\[
 \operatorname{Ext}^q_{\operatorname{Gr}D}(S_i,P_j)=0
 \quad\text{unless }q=3,\ j=i+4,
 \qquad
 \operatorname{Ext}^3(S_i,P_{i+4})\simeq W_i^\vee.
\]
最後の同定は解消の最高項への Hom から得られる。
各固定次数ごとの議論で十分であり、無限個の次数を同時に自明化する
一つの開被覆は必要ない。
\end{proof}

普遍 point curve を $\pi:C\to S$ とする。
隣接二因子の評価による $2\times2$ 行列の determinant が
$\PP(V_0^\vee)\times\PP(V_1^\vee)$ 内に $C$ を定める。
geometric 条件は、determinant が零である点でもこの行列の階数が一であることを保証する。
三点 correspondence の隣接二点への射影は、単元小行列式による掃き出しから
スキームとして同型になる。
従って全 $\mathcal L_i$ を同じ $C$ 上の直線束として定義できる。
これは第~\ref{sec:relative-cubic-noetherian}節の相対 point-curve 構成である。

$S$ は線型 locus を除いているので、determinant は各幾何ファイバーで非零である。
したがって $C/S$ は平坦な relative Cartier divisor である。
その標準 exact sequence
\[
 0\longrightarrow\OO_{\PP^1_S\times\PP^1_S}(-2,-2)\otimes\pi^*\ell
 \longrightarrow\OO_{\PP^1_S\times\PP^1_S}
 \longrightarrow\OO_C\longrightarrow0
 \tag{PR-3}\label{eq:curve-sequence}
\]
（$\ell$ は基底の可逆層）から、$\pi_*\OO_C=\OO_S$ が普遍的に成り立ち、
$R^1\pi_*\OO_C$ は可逆層、$R^{>1}\pi_*\OO_C=0$ となる。

\begin{lemma}[被約基底での相対 Picard の一致判定]\label{lem:reduced-equalizer}
$X$ を被約スキーム、$Y\to X$ を相対 diagonal が immersion である代数空間とする。
二つの section $f,g:X\to Y$ が全幾何点上で一致すれば、$f=g$ である。
この判定は $Y=\Pic_{C/X}$ に適用できる。
\end{lemma}

\begin{proof}
equalizer $E\to X$ は diagonal の底変更なので immersion である。
仮定よりその像は $|X|$ 全体である。
従って局所的には $X$ の全ての点で零となる ideal による閉 immersion であり、
$X$ の被約性からその ideal は零である。
相対 Picard の diagonal は quasi-compact immersion である
\cite[Lemma 108.9.2]{StacksRelativePicard}。
相対 Picard 全体の separatedness は必要ない。
\end{proof}

\begin{lemma}[普遍 helix の相対的な差]\label{lem:relative-helix-difference}
$N=[\mathcal L_2\mathcal L_0^{-1}]$ は $P=\Pic^0_{C/S}$ の section であり、相対 Picard 群で
\[
 [\mathcal L_{i+2}\mathcal L_i^{-1}]=N
 \tag{PR-4}\label{eq:helix-difference}
\]
が全ての $i$ について成り立つ。
\end{lemma}

\begin{proof}
各幾何ファイバーの admissibility により、$\mathcal L_2$ と $\mathcal L_0$ の多重次数は等しい。
従って $N$ は $P$ に属する。
体上の elliptic helix は
$\mathcal L_i\mathcal L_{i+1}^{-1}\mathcal L_{i+2}^{-1}\mathcal L_{i+3}\simeq\OO_C$
を満たすので、隣接する二つの差
$[\mathcal L_{i+2}\mathcal L_i^{-1}]$ と $[\mathcal L_{i+3}\mathcal L_{i+1}^{-1}]$ は等しい。
全ての差は $N$ に等しく、被約な $S$ に補題~\ref{lem:reduced-equalizer} を適用して
\eqref{eq:helix-difference} を相対的に得る。
有限個の $i$ についてこの議論を行えば、以下の二周期の構成には十分である。
\end{proof}

\subsection{平方根の étale パラメータ空間}

\begin{lemma}[相対 Picard の二等分]\label{lem:picard-halves}
$P=\Pic^0_{C/S}$ は滑らかな相対次元一の可換群代数空間であり、
\[
 [-2]:P\longrightarrow P
\]
は étale かつ全射である。
従って
\[
 T=S\mathbin{\times}_{N,P,[-2]}P
 \tag{PR-5}\label{eq:T}
\]
は $S$ 上 étale 全射であり、普遍 section $a\in P(T)$ は
$a^{-2}=N$、すなわち $a^2=N^{-1}$ を満たす。
\end{lemma}

\begin{proof}
\eqref{eq:curve-sequence} と曲線の相対 Picard の一般定理により
$\Pic_{C/S}$ は代数空間で局所有限表示であり、相対次元一の曲線であるため滑らかである
\cite[Lemma 108.9.5]{StacksRelativePicard}。
相対 Picard は quasi-separated でもあるため、
\cite[Appendix A, Lemma A.1]{HolmesMolchoOrecchiaPoiret2023} を適用できる。
その fiberwise identity component $P$ は一意な開部分群である。
開部分と幾何ファイバーの恒等成分による特徴付けの一意性から、任意の底変更と両立する。
各ファイバーの接空間は $H^1(C_s,\OO_{C_s})$ なので相対次元は一である。

ここで多重次数零と identity component の一致は非被約ファイバーも含む。
体上で nilradical を $I\subset\OO_C$ とすると、標数零の有限な対数級数により
$1+I\simeq I$ が成り立つ。
$H^2(C,I)=0$ なので $\Pic C\to\Pic C_{\mathrm{red}}$ は全射であり、
その核は $H^1(C,I)$ から来る連結な加法群である。
被約曲線では正規化と conductor の単元の完全列により、
多重次数零の部分は、正規化の各成分の Jacobian の積の
連結な torus・unipotent 群による拡大である。
従って多重次数零の部分はちょうど連結成分 $\Pic^0(C)$ となる。
相対的には滑らかな群代数空間の fiberwise identity component の開部分群を取っている。

二倍写像の相対微分は群の単位元で $2\,\mathrm{id}$、他の点ではその平行移動である。
$2$ は可逆なので、滑らかな相対一次元空間の間の $[-2]$ は étale となる。
各幾何ファイバーでその像は開部分群である。
ファイバー $P_s$ は連結なので、この像は全体である。
従って $[-2]$ は全射である。
\eqref{eq:T} の主張は底変更から従う。
\end{proof}

$T$ は $k$ 上滑らかで被約である。
必要なら $T$ の étale scheme chart を取り、以下では $T$ 自体をスキームとして扱う。
命題~\ref{prop:ATV-relative-action} により、
普遍 Picard class $a$ から実際の曲線の自己同型
\[
 \sigma=\eta(a):C_T\longrightarrow C_T
 \tag{PR-6}\label{eq:sigma}
\]
が得られる。
ここでは $a$ を表す大域的な直線束を選ぶ必要はない。
群作用そのものが相対 Picard 空間上に定義されている。

\begin{lemma}[線束の比較と二周期]\label{lem:relative-beta}
$T$ の Zariski 開被覆上で同型
\[
 \varepsilon_i:\sigma^*\mathcal L_i\xrightarrow{\sim}\mathcal L_{i+2}
 \qquad(i=0,1,2,3)
 \tag{PR-7}\label{eq:epsilons}
\]
を選べ、それらが実際の代数同型
\[
 \beta:D_T\xrightarrow{\sim}D_T(2)
 \tag{PR-8}\label{eq:beta}
\]
を与える。
\end{lemma}

\begin{proof}
全幾何点で \eqref{eq:field-eta} と $\deg\mathcal L_i=2$ により
\[
 [\sigma^*\mathcal L_i]=[\mathcal L_i]a^{-2}=[\mathcal L_i]N=[\mathcal L_{i+2}]
\]
である。
被約な $T$ 上で補題~\ref{lem:reduced-equalizer} を適用すると、これは
相対 Picard 空間の section としての等式である。

従って $\sigma^*\mathcal L_i\otimes\mathcal L_{i+2}^{-1}$ は基底の直線束の引戻しである。
詳しくは、相対 Picard class が零なので fppf 局所に自明となる。
$\pi_*\OO_{C_T}=\OO_T$ が普遍的に成り立つことから、その自明化の変換関数は
基底から来るため、直線束そのものが基底へ降下する。
四つの比較直線束を Zariski 局所に自明化して\eqref{eq:epsilons} を得る。

$\pi_*\mathcal L_i$ は $V_i\otimes\OO_T$ と同一視される。
これは全ファイバーで $H^1(\mathcal L_i)=0$、$h^0(\mathcal L_i)=2$ であり、
普遍 evaluation との比較がファイバー上同型になることから従う。
\eqref{eq:epsilons} による切断の引戻しは可逆写像
$\beta_i:V_i\to V_{i+2}$ を与える。
三次関係は
\[
 R_i=\ker\left(V_i\otimes V_{i+1}\otimes V_{i+2}
   \longrightarrow\pi_*(\mathcal L_i\mathcal L_{i+1}\mathcal L_{i+2})\right)
\]
であるから、$\beta_0,\ldots,\beta_3$ は $R_0,R_1$ を $R_2,R_3$ へ移す。
従ってその tensor 積は $W_0$ を $W_2$ へ移す。
四周期により $\beta_{i+4}=\beta_i$ と延長する。
巡回を $c$ と書けば、対応する tensor 積の射について
\[
 c(\beta_0\otimes\beta_1\otimes\beta_2\otimes\beta_3)
  =(\beta_1\otimes\beta_2\otimes\beta_3\otimes\beta_0)c
\]
が成り立つ。
従って $W_0$ を $W_2$ へ移すことから、巡回された全ての $W_i$、
従って全ての三次関係を保つ。
逆写像も同じ性質を持ち、\eqref{eq:beta} を得る。
最初の隣接 pair に対する曲線の作用は、構成から $\sigma$ の作用である。
三点 correspondence の二つの隣接 pair への射影はスキーム同型であり、
代数同型による作用と可換する。
この自然性により、全ての添字で同じ $\sigma$ が誘導される。
従って四周期に延長した $\beta$ の平方も、切断上の $\sigma^{2*}$ と projectively 一致する。
\end{proof}

\subsection{平方と canonical 四周期の比較}

\begin{theorem}[相対的な適合 Serre 根]\label{thm:relative-compatible-root}
上の $T$ をさらに Zariski 局所化した基底上で、
\eqref{eq:beta} は頂点対角型の自己同型 $\delta$ に対し
\[
 \beta^2=\delta\theta
 \tag{PR-9}\label{eq:beta-square}
\]
を満たす。
従って tails 上では $(2)_\beta^2\simeq(4)_\theta$ である。
この同型は基底上に構成されており、任意の底変更後にも成り立つ。
特に非被約基底へ引き戻しても、相対 Serre 関手を持つ有限 foundation の比較の下で
\[
 s=(2)_\beta,\qquad s^2\simeq S^{-1}[2]
 \tag{PR-10}\label{eq:Serre-square}
\]
を得る。
\end{theorem}

\begin{proof}
\emph{(a) 四周期と Picard の符号.}
$\eta$ は可換群の作用なので、$T$ 上で厳密に
\[
 \sigma^2=\eta(a)^2=\eta(a^2)=\eta(N^{-1})
 \tag{PR-11}\label{eq:sigma-square}
\]
である。
ただしこの式だけで canonical $\theta$ との同定を済ませない。
その同定には次の tensor の比較を使う。

\emph{(b) smooth elliptic locus での tensor 比較.}
$S_{\mathrm{sm}}\subset S$ を $C/S$ が滑らかな曲線の locus とする。
これは非空開部分で、既約な $S$ において稠密である。
その上の各幾何点では、補題~\cite[\texttt{lem:elliptic-period-square}]{QuadricIntrinsic2026} の次の計算を使う。
原点を選び、$\mathcal L_i$ の次数二の因子和を $b_i$ と書くと
$b_{i+2}-b_i$ は一定である。
$\sigma$ は $2v=b_i-b_{i+2}$ を満たす平行移動 $p\mapsto p+v$ である。
$\mathcal L_{i+2}$ の二重被覆の involution を
$\iota_{i+2}(p)=b_{i+2}-p$ とする。

$w_i$ を一般点 $p$ の $\mathcal L_{i+2}$ の評価で縮約し、
$p'=\iota_{i+2}(p)$ を使って二つの rank-one summand に分ける。
隣接二因子の縮約の非零性により、その両係数は非零である。
$w_{i+2}$ が $w_i$ の二回巡回であることから、
$\theta_i$ による $\mathcal L_{i+4}$ の $p$ における評価の引戻しは
$\mathcal L_i$ の $p'$ における評価である。
後者は $\mathcal L_i$ の involution を使うと
\[
 \iota_i(p')=p+b_i-b_{i+2}=p+2v
\]
における評価と同じである。
従って $\theta_i$ と $\beta_{i+2}\beta_i$ は $V_i$ 上で projectively 等しい。
この比較は切断の引戻しという \eqref{eq:epsilons} の向きと一致する。

\emph{(c) 普遍基底上での延長.}
$T\to S$ は étale なので、$T_{\mathrm{sm}}$ は $T$ の各既約成分で稠密である。
実際、$T$ は滑らかであり、その各成分の像は $S$ の非空開部分を含む。
四つの projective linear map
\[
 [\beta_{i+2}\beta_i],\ [\theta_i]
 \quad(i\bmod4)
\]
は $\PGL(V_i,V_{i+4})$ という separated な空間への射である。
(b) により $T_{\mathrm{sm}}$ の全幾何点で一致し、
$T$ は被約なので、全 $T$ で一致する。
ここで被約なのは普遍根の基底 $T$ である。
任意の非被約な底変更先で、幾何ファイバーの一致だけから結論しているのではない。

\emph{(d) 頂点スカラーの明示.}
したがって一意な単元 $c_i\in\OO_T^\times$ があり、
\[
 \beta_{i+2}\beta_i=c_i\theta_i,
 \qquad c_{i+4}=c_i
\]
となる。
$u_0=1$、$u_{i+1}=c_i^{-1}u_i$ と全整数について再帰的に定める。
各 $u_i$ は有限積からなる単元である。
すると
\[
 \delta|_{D_{ij}}=\frac{u_i}{u_j}\,\mathrm{id}_{D_{ij}}
\]
は代数自己同型で、一次生成性により \eqref{eq:beta-square} が成り立つ。
右加群の $i$ 成分を $u_i$ 倍する射は、$\delta$ による引戻しと恒等関手の自然同型を与える。
従って tails 上の平方同型を得る。

\emph{(e) 基底上の Serre 双加群比較.}
定理~\ref{thm:relative-cubic-tails-tilting} により、連続四項の直和は
相対 compact tilting generator となる。
その有限 foundation 代数を $\Lambda$ とし、
$s^2$ と相対反 Serre 関手 $S^{-1}[2]$ を表す
可逆 perfect $\Lambda$-双加群を $P,Q$ とする。
$P$ は実際の次数移動 $s=(2)_\beta$ を標準 dg Morita 同値で移したものである。
$Q$ は $\Lambda$ の相対 Serre 双加群の逆と $[2]$ から得られる。
有限 foundation の相対 smooth proper 性
（補題~\ref{lem:foundation-relative-smooth-proper}）により、
これらは基底変更と整合する可逆 perfect 双加群である。

全ての幾何点で (a)--(d) が成り立つので、
\emph{ここで構成した同じ $\beta$} に
命題~\cite[\texttt{prop:admissible-period-h2}]{QuadricIntrinsic2026} を適用できる。
従って $P_{\bar t}\simeq Q_{\bar t}$ である。
これは各ファイバーに何らかの根が存在するという主張より強く、
普遍次数移動 $s$ の平方そのものの比較である。

任意の底変更後にも $Z(\Lambda)=\mathcal O_T$ である。
実際、中心元は四つの頂点で scalar であり、
隣接する忠実な有限局所自由 Hom 加群との可換性により全 scalar が等しい。
従って定理~\ref{thm:bimodule-line-comparison} を $P,Q$ に適用できる。
比較 locus $T_{P\sim Q}\hookrightarrow T$ は有限表示局所閉 immersion で、
全ての幾何点を含む。
$T$ は被約なので、補題~\ref{lem:reduced-equalizer} の proof と同じ ideal の議論により
$T_{P\sim Q}=T$ となる。
この locus の構成は、最低次 Hom の行列を基底環で零にする条件を含むので、
非被約基底に対しても有効な scheme-theoretic 比較である。

したがって比較線束
\[
 \ell=H^0\mathbf R\mathcal Hom_{\Lambda^e}(P,Q)
\]
と評価同型 $P\otimes\ell\simeq Q$ を得る。
さらに Zariski 局所に $\ell$ を自明化すれば
\eqref{eq:Serre-square} が成り立つ。
全ての射は普遍根の基底上に構成されているため、
任意の非被約基底への引戻し後も同じ同型が成り立つ。
\end{proof}



\begin{remark}[比較線束の局所自明化]
相対 Picard の等式からは、まず二つの直線束の差が基底から来ることが分かる。
\eqref{eq:epsilons} の前に四つの比較直線束を自明化したことが必要である。
その選択を変えると $\beta_i$ は基底単元倍されるが、
projective な平方比較は変わらず、(d) の頂点スカラーがその差を吸収する。
さらに (e) では Serre 双加群の比較線束 $\ell$ を自明化している。
従って幾何ファイバー上の平方根の存在だけから、
基底線束を省略した大域的な平方同型を主張してはいない。
必要なのは有限個の基底直線束の局所自明化である。
\end{remark}

\begin{remark}[四項 foundation と右加群の向き]
Abelian helix から $D_{ij}=\Hom(E_i,E_j)$ を回復する場合は、
右 tails の構成に添字を反転した反対代数
\[
 C_{ij}=D_{-j,-i},\qquad \beta_C=\beta_D^{-1}
\]
を用いる。
すると $E_i$ は $\pi e_{-i}C$ に対応し、四項窓の自己準同型代数は
元の foundation の積と一致する。
$\beta_D^2=\delta\theta_D$ を反転すれば、$C$ 上でも
$\beta_C^2=\theta_C$ が頂点スカラーを除いて成り立つ。
ここで $\theta_C$ は $\theta_D^{-1}$ に対応する canonical 四周期である。
(e) の双加群 $P,Q$ とその比較 locus は、この $C$ の右 tails 上で直接取れる。
従って Serre の比較と右軌道積の規約は同じ向きで成立する。
\end{remark}

\subsection{Artin 持上げと atlas 構成への入力}

\begin{corollary}[有限型の étale 根 chart]\label{cor:finite-type-root-chart}
$S$ を覆う有限個の affine scheme $T_\alpha$ と étale 射 $T_\alpha\to S$ を選び、
その各々で上の実際の二周期および Serre 平方同型を構成できる。
従って $\coprod_\alpha T_\alpha$ は有限型の滑らかな $k$ スキームであり、
全ての非線型 cubic geometric 五つ組の点を覆う。
\end{corollary}

\begin{proof}
$T\to S$ は étale 全射であり、その scheme chart の affine open の像は
$S$ を開被覆する。$S$ は準コンパクトなので有限個を選べる。
各 affine chart は $S$ 上有限型に取れる。
四つの Picard 比較直線束と Serre 比較線束 $\ell$ を
有限個の affine open で自明化し、必要な被覆を細分する。
\end{proof}

\begin{corollary}[指定した Picard 根の全 Artin 持上げ]\label{cor:all-artin-root}
$x\in S(k)$ と $a_x\in\Pic^0(C_x)$、$a_x^{-2}=N_x$ を固定する。
\eqref{eq:T} の中の $a_x$ に対応する点を含む affine étale scheme chart を選び、
その点を $t$ とする。必要なら前の有限被覆にこの chart を一枚追加する。
任意の局所 Artin $k$-代数 $A$ と、その閉点が $x$ となる framed geometric
五つ組 $\operatorname{Spec} A\to S$ に対し、$t$ を指定した $\operatorname{Spec} A\to T$ の持上げは一意に存在する。
対応する二周期は Serre 適合であり、任意の Artin 商・小拡大に整合する。
\end{corollary}

\begin{proof}
étale 射は nilpotent closed immersion に対して形式的 étale である。
閉点の持上げ $t$ から $\operatorname{Spec} A\to T$ が一意に得られる。
さらに選んだ Zariski 近傍は $t$ を含むので、$\operatorname{Spec} A$ 全体がその近傍へ入る。
そこで構成した \eqref{eq:beta-square} と \eqref{eq:Serre-square} を引き戻す。
\end{proof}

\paragraph{Abelian 変形と atlas への適用.}
相対根の構成自体は、非線型 locus 全体で特異・可約・非被約な曲線ファイバーを含む。
平坦 Abelian Artin 変形からの適用では、
補題~\cite[\texttt{lem:cubic-artin-overlap-recovery}]{QuadricIntrinsic2026} と命題~\cite[\texttt{prop:framed-cubic-artin-lift}]{QuadricIntrinsic2026} が
DDL による helix の持上げから、可逆部分束 $W_0$ とその overlap・縮約を含む
指定基底付き五つ組を回復する。
補題~\ref{lem:artin-two-heart-comparison} は、同じ四項 foundation を持つ二つの圏の
標準 dg Morita 同値が、指定された exact な還元を保って heart-exact となることを示す。
従って系~\ref{cor:all-artin-root} の根を元の Abelian 変形へ移せる。
この指定同型の持上げと有限型効果性を合わせた smooth atlas は
定理~\ref{thm:quintuple-chart-artin-lifting} と定理~\ref{thm:cubic-smooth-open} で構成する。
線型 cubic 点は全て pointed $(\operatorname{coh}(\PP^1\times\PP^1),\OO)$ に対応し、
既存の可換二次曲面の étale Morita-root chart を用いればよい。

